% Generated draft: requires mathematical and source review.
\documentclass[11pt]{article}
\usepackage[margin=25mm]{geometry}
\usepackage{fontspec}
\setmainfont{Latin Modern Roman}
\usepackage{amsmath,amssymb,mathtools,unicode-math}
\setmathfont{Latin Modern Math}
\usepackage{enumitem,needspace,xurl,hyperref}
\setlength{\parindent}{0pt}
\setlength{\parskip}{5pt}
\setlength{\emergencystretch}{3em}
\newcommand{\N}{\mathbb{N}}
\newcommand{\Z}{\mathbb{Z}}
\newcommand{\Q}{\mathbb{Q}}
\newcommand{\R}{\mathbb{R}}
\newcommand{\C}{\mathbb{C}}
\newcommand{\F}{\mathbb{F}}
\newcommand{\A}{\mathbb{A}}
\newcommand{\PP}{\mathbb{P}}
\newcommand{\E}{\mathbb{E}}
\newcommand{\eps}{\varepsilon}
\newcommand{\dd}{\,\mathrm{d}}
\begin{document}
\section*{UnsolvedMath ID: 30005685}
\textbf{Source problem code:} catalogue entry
\subsection{Definitions and mathematical statement}
The source identifies the Fibonacci Hamiltonian as a discrete Schr\"odinger operator whose potential is generated by the Fibonacci substitution sequence. It does not specify the lattice, boundary conditions, potential values, coupling parameters, or the law and operator form of the random perturbation \(W_\omega\). Consequently, the supplied record does not determine a unique random operator model.

A conditional formulation of the intended spectral question is as follows. Fix the intended self-adjoint realization \(H\) on a Hilbert space \(\mathcal H\), a probability space \((\Omega,\mathcal F,\mathbb P)\), and a random perturbation \(W_\omega\), including operator domains, such that \(A_\omega=H+W_\omega\) is a well-defined self-adjoint operator. For a self-adjoint operator \(A\), write \(\sigma(A)\) for its spectrum. If there is a deterministic closed set \(\Sigma\subseteq\R\) satisfying
\[
\sigma(A_\omega)=\Sigma
\qquad\text{for \(\mathbb P\)-almost every \(\omega\in\Omega\)},
\]
then \(\Sigma\) is its almost-sure spectrum. The requested problem is to determine the topological structure of \(\Sigma\) for the Fibonacci Hamiltonian and random perturbation intended in the source. This is a conditional description of the target, not an assertion that every random perturbation admits a deterministic almost-sure spectrum. The missing model specifications must be recovered before the question becomes a precise model-specific problem.

The supplied excerpt also states, as Corollary 3, that if a base space \(X\) is connected, then an almost-sure spectrum denoted there by \(\Sigma_1\) is a finite union of non-degenerate compact intervals. Here a non-degenerate compact interval means \([a,b]\) with \(a<b\). The excerpt does not supply the definition of \(X\), the random model, or the other hypotheses of that corollary. Its conclusion therefore cannot be asserted here from connectedness alone. The supplied literature assessment says that connected-base theorems do not apply to the Fibonacci hull; this is a supplied qualification, not an independently verified theorem application.
\subsection{Short English statement}
Determine the topology of the almost-sure spectrum of a randomly perturbed Fibonacci Hamiltonian. The supplied record labels this question open, but omits the perturbation law and several operator specifications needed for a precise formulation. The connected-base interval conclusion quoted nearby cannot be transferred to this problem using the supplied excerpt alone.
\subsection{Sources}
\begin{itemize}
\item Local UnsolvedMath registry record, permanent ID \(30005685\), source problem code \texttt{catalogue entry}, title \emph{Spectrum of Randomly Perturbed Fibonacci Hamiltonians}. Project registry paths: \texttt{data/unsolvedmath/problems.json} and supplementary \texttt{catalog/problems.json}. This draft uses the supplied complete record; registry files were not independently inspected.
\item Supplied primary-source citation: \emph{Aspects of Aperiodic Order (2024)}, Oberwolfach Reports, DOI \texttt{10.4171/owr/2023/37}, pp.\ 2107--2178, \url{https://doi.org/10.4171/owr/2023/37}. The supplied original excerpt identifies Question 4 and quotes an adjacent Corollary 3. No individual proposer is supplied. The record gives proposed year 2023, while its report citation uses 2024; both are retained without reconciliation.
\item Supplied literature source labeled Avila--Damanik--Gorodetski: \url{https://arxiv.org/abs/2211.02173}.
\item Supplied literature source labeled Damanik--Gorodetski--Yessen: \url{https://arxiv.org/abs/1403.7823}.
\item Literature assessment date: \texttt{2026-08-22}. This is supplied dated curation, not independently verified. It labels the question open, states that connected-base theorems do not apply to the Fibonacci hull, and reports that no primary result was found giving the full spectral topology. The supplied background explicitly requests independent verification of time-sensitive status and cited claims.
\item The registry marks the statement \texttt{corrected\_verified} and says that it was reconstructed and checked against the official report. These are supplied provenance claims; this draft does not independently verify them or the mathematical open status.
\end{itemize}
\end{document}
